% Define Sans Serif para todo documento
\renewcommand{\familydefault}{\sfdefault}

% Deve haver recuo no primeiro paragrafo
\usepackage{indentfirst}

% Capitulos, secoes e subsecoes devem ter tamanho normal
\renewcommand{\ABNTEXchapterfontsize}{\normalsize}
\renewcommand{\ABNTEXsectionfontsize}{\normalsize}
\renewcommand{\ABNTEXsubsectionfontsize}{\normalsize}
\renewcommand{\ABNTEXsubsubsectionfontsize}{\normalsize}

% Capitulos e secoes devem estar negritados
% https://groups.google.com/g/latex-br/c/ULm17Zfng3w/m/ZbCHIHLi7GAJ
\renewcommand{\ABNTEXchapterfont}{\normalfont\sffamily\bfseries}
\renewcommand{\ABNTEXsectionfont}{\normalfont\sffamily\bfseries}
% Subsecoes nao devem estar negritados
\renewcommand{\ABNTEXsubsectionfontsize}{\normalfont\sffamily}

% Espacamento vertical entre titulo do capitulo e conteudo deve ser de 24pts
\setlength\afterchapskip{24pt}

% Define cores para links
\hypersetup{
    colorlinks=true,
    linkcolor=black,
    citecolor=black,
    urlcolor=black,
}

% Permite quebra das URLs após qualquer caractere (evita estouro de margem).
\usepackage{xurl}

% URLs sublinhadas. O \underLine do lua-ul (requer LuaLaTeX) desenha um
% sublinhado tipográfico que acompanha as quebras de linha da URL em qualquer
% ponto, sem encaixotar o conteúdo (ao contrário do ulem/soul). \hyper@normalise
% sanitiza a URL (trata & ? # etc.) antes de processá-la.
%
% IMPORTANTE: lua-ul/luacolor só funcionam com LuaLaTeX. Sob pdfLaTeX eles
% travam a compilação (\LuaULResetUnderline fica indefinido). Por isso só
% carregamos este bloco quando o motor for LuaTeX; sob pdfLaTeX as URLs apenas
% não ficam sublinhadas, mas o documento compila normalmente.
\ifdefined\directlua
  \usepackage{luacolor}
  \usepackage{lua-ul}
  \makeatletter
  \let\orig@url\url
  \renewcommand{\url}{\hyper@normalise\ul@url@aux}
  \newcommand{\ul@url@aux}[1]{\underLine{\hyper@linkurl{{\UrlFont\nolinkurl{#1}}}{#1}}}
  \makeatother
\fi

% Páginas sem cabeçalho
\pagestyle{simple}

% Profundidade do sumário: 3 inclui seções, subseções e subsubseções
\addtocontents{toc}{\protect\setcounter{tocdepth}{3}}

% Recuo à esquerda dos subitens e subsubitens no Sumário.
% (indent = deslocamento da entrada; numwidth = largura reservada ao número)
\setlength{\cftsubsectionindent}{1.8em}
\setlength{\cftsubsectionnumwidth}{3.3em}
\setlength{\cftsubsubsectionindent}{4.2em}
\setlength{\cftsubsubsectionnumwidth}{4.2em}

% Captions com negrito
% \usepackage[font=bf]{caption}
% Caption pequena, com o "Figura" em negrito
\usepackage[font=small,labelfont=bf]{caption}

% Posicionamento de Figuras - Importantíssimo
\usepackage{float}

% Espaçamento de 1,5
\renewcommand{\baselinestretch}{1.5}
